% NOTE (yours to reframe): this is a fact-correct scaffold, not final framing.
% The goal sentence and motivation are where your taste sits -- rewrite in your voice.
\section{Introduction}
Post-quantum anonymous credentials let a holder prove possession of a valid
credential without revealing identity, under assumptions that survive a
quantum adversary. We target the setting in which the holder generates this
proof on a \emph{personal device}, here a consumer laptop (18-core, 24\,GB),
as opposed to a cloud or server prover. A general-purpose zero-knowledge
virtual machine (zkVM) is attractive in this setting for \emph{flexibility}
rather than for raw speed: it proves an arbitrary verification predicate without
a bespoke, per-scheme circuit. The obstacle is arithmetic. The schemes we
care about (PLUM, Loquat) operate over a $192$-bit prime field, while
production zkVMs (RISC~Zero, SP1) fix a small ($\approx 31$-bit) native
field, so every operation is emulated by multi-limb arithmetic; an algebraic
hash (Griffin) accounting for $\approx 91\%$ of the verification cost
amplifies the penalty. This thesis asks whether such a verification can be
made to run on a personal device at all, what it costs, and whether the
resulting proof can additionally be zero-knowledge.
